% Options for packages loaded elsewhere
\PassOptionsToPackage{unicode}{hyperref}
\PassOptionsToPackage{hyphens}{url}
\documentclass[
  english,
  11pt,
]{article}
\usepackage{xcolor}
\usepackage[margin=1in]{geometry}
\usepackage{amsmath,amssymb}
\setcounter{secnumdepth}{-\maxdimen} % remove section numbering
\usepackage{iftex}
\ifPDFTeX
  \usepackage[T1]{fontenc}
  \usepackage[utf8]{inputenc}
  \usepackage{textcomp} % provide euro and other symbols
\else % if luatex or xetex
  \usepackage{unicode-math} % this also loads fontspec
  \defaultfontfeatures{Scale=MatchLowercase}
  \defaultfontfeatures[\rmfamily]{Ligatures=TeX,Scale=1}
\fi
\usepackage{lmodern}
\ifPDFTeX\else
  % xetex/luatex font selection
\fi
% Use upquote if available, for straight quotes in verbatim environments
\IfFileExists{upquote.sty}{\usepackage{upquote}}{}
\IfFileExists{microtype.sty}{% use microtype if available
  \usepackage[]{microtype}
  \UseMicrotypeSet[protrusion]{basicmath} % disable protrusion for tt fonts
}{}
\makeatletter
\@ifundefined{KOMAClassName}{% if non-KOMA class
  \IfFileExists{parskip.sty}{%
    \usepackage{parskip}
  }{% else
    \setlength{\parindent}{0pt}
    \setlength{\parskip}{6pt plus 2pt minus 1pt}}
}{% if KOMA class
  \KOMAoptions{parskip=half}}
\makeatother
\usepackage{longtable,booktabs,array}
\newcounter{none} % for unnumbered tables
\usepackage{calc} % for calculating minipage widths
% Correct order of tables after \paragraph or \subparagraph
\usepackage{etoolbox}
\makeatletter
\patchcmd\longtable{\par}{\if@noskipsec\mbox{}\fi\par}{}{}
\makeatother
% Allow footnotes in longtable head/foot
\IfFileExists{footnotehyper.sty}{\usepackage{footnotehyper}}{\usepackage{footnote}}
\makesavenoteenv{longtable}
\usepackage{graphicx}
\makeatletter
\newsavebox\pandoc@box
\newcommand*\pandocbounded[1]{% scales image to fit in text height/width
  \sbox\pandoc@box{#1}%
  \Gscale@div\@tempa{\textheight}{\dimexpr\ht\pandoc@box+\dp\pandoc@box\relax}%
  \Gscale@div\@tempb{\linewidth}{\wd\pandoc@box}%
  \ifdim\@tempb\p@<\@tempa\p@\let\@tempa\@tempb\fi% select the smaller of both
  \ifdim\@tempa\p@<\p@\scalebox{\@tempa}{\usebox\pandoc@box}%
  \else\usebox{\pandoc@box}%
  \fi%
}
% Set default figure placement to htbp
\def\fps@figure{htbp}
\makeatother
\ifLuaTeX
\usepackage[bidi=basic,shorthands=off]{babel}
\else
\usepackage[bidi=default,shorthands=off]{babel}
\fi
\ifLuaTeX
  \usepackage{selnolig} % disable illegal ligatures
\fi
\setlength{\emergencystretch}{3em} % prevent overfull lines
\providecommand{\tightlist}{%
  \setlength{\itemsep}{0pt}\setlength{\parskip}{0pt}}
% Map Unicode symbols that the text font lacks to proper math glyphs
\usepackage{newunicodechar}
\newunicodechar{≥}{\ensuremath{\geq}}
\newunicodechar{≤}{\ensuremath{\leq}}
\newunicodechar{≈}{\ensuremath{\approx}}
\newunicodechar{χ}{\ensuremath{\chi}}
\newunicodechar{±}{\ensuremath{\pm}}
\newunicodechar{−}{\ensuremath{-}}
\newunicodechar{≠}{\ensuremath{\neq}}
% Keep wide baseline/results tables within the margins
\usepackage{etoolbox}
\AtBeginEnvironment{longtable}{\small}
\usepackage{bookmark}
\IfFileExists{xurl.sty}{\usepackage{xurl}}{} % add URL line breaks if available
\urlstyle{same}
\hypersetup{
  pdftitle={Machine-learning estimation of obstructive coronary artery disease from routine clinical data --- a development and internal-validation study},
  pdflang={en},
  hidelinks,
  pdfcreator={LaTeX via pandoc}}

\title{Machine-learning estimation of obstructive coronary artery
disease from routine clinical data --- a development and
internal-validation study}
\author{}
\date{}

\begin{document}
\maketitle

Articles · Cardiology · Clinical prediction model

\section{Machine-learning estimation of obstructive coronary artery
disease from routine clinical and laboratory data to support
pre-angiographic triage: development and internal validation of the
CardioPredict
model}\label{machine-learning-estimation-of-obstructive-coronary-artery-disease-from-routine-clinical-and-laboratory-data-to-support-pre-angiographic-triage-development-and-internal-validation-of-the-cardiopredict-model}

{[}Author One{]}\textsuperscript{a}, {[}Author
Two{]}\textsuperscript{b}, {[}Author Three{]}\textsuperscript{a},
{[}Senior Author{]}\textsuperscript{a,*}

\textsuperscript{a} {[}Department / Institution{]}, {[}City{]},
{[}Country{]}. ~ \textsuperscript{b} {[}Second Institution{]},
{[}City{]}, {[}Country{]}.

\textsuperscript{*} Correspondence: {[}Senior Author{]}, {[}address{]},
{[}email{]}.

Analysis seed fixed (42) · Reported in accordance with the TRIPOD
statement\textsuperscript{10} for multivariable prediction models. Development cohort data
hash (SHA-256, first 16 hex): {fb3af1d8645920f3}.

\subsection{Summary}\label{summary}

\textbf{Background} Invasive coronary angiography is the reference
standard for obstructive coronary artery disease (CAD) but is costly,
carries procedural risk, and a substantial proportion of diagnostic
angiograms reveal no obstructive disease. Tools that quantify the
pre-test probability of obstructive CAD from data available before
catheterisation could improve patient selection. We developed and
internally validated a machine-learning model (CardioPredict) that
estimates the probability of obstructive CAD using only routine
demographic, clinical and laboratory variables.

\textbf{Methods} In this single-centre retrospective study we analysed
1187 patients referred for coronary angiography (the analytic sample,
drawn from 2346 registry entries after excluding those without an
angiographic outcome). The outcome was angiographically obstructive CAD.
Thirty pre-procedural variables (demographics, anthropometry, blood
pressure, lipid profile, glycaemic indices, inflammatory and
haematological markers, and personal/family history) were used; all
angiography-derived and outcome-derived fields were removed before
modelling to prevent label leakage. A four-learner ensemble (TabPFN,
gradient-boosted trees {[}LightGBM and CatBoost{]}, and a clinical
logistic-regression model) was combined through an age-conditional
stacked meta-learner, with stratified probability calibration and
split-conformal prediction for case-level uncertainty. Discrimination,
calibration and clinical-utility metrics were estimated by 2 repeats of
5-fold nested cross-validation, which re-fits all preprocessing and
learners within each training fold, with 95\% CIs obtained by bootstrap.
We benchmarked the ensemble against a penalised logistic-regression
model fitted on the same predictors and folds. The deployed model was
then fitted on the full cohort with all base learners persisted for
deterministic inference.

\textbf{Findings} Obstructive CAD was present in 781 of 1187 patients
(65·8\%). Using 30 routine pre-procedural predictors and repeated nested
cross-validation with fold-internal imputation, the model achieved an
area under the receiver-operating-characteristic curve (AUC) of 0·866
(SD across repeats 0·001; bootstrap 95\% CI 0·85--0·89). At the
prespecified operating threshold (0·73, maximising the Matthews
correlation coefficient {[}MCC{]}), sensitivity was 74·0\% (95\% CI
70·8--77·0), specificity 81·5\% (77·6--85·2), positive predictive value
88·5\% (86·0--91·0), negative predictive value 62·0\% (58·0--66·0),
overall accuracy 76·6\% (74·1--79·0), and MCC 0·53. Predicted
probabilities were well calibrated (calibration slope 0·91;
calibration-in-the-large −0·05; expected/observed {[}E/O{]} ratio 1·02).
A penalised logistic-regression model on the same engineered features
reached an AUC of 0·794 (0·77--0·82); the ensemble\textquotesingle s
incremental AUC was +0·077 (95\% CI +0·059 to +0·098; p\textless0·001).
The conformal framework flagged 23·7\% of patients as statistically
indeterminate (gray-zone), with an empirical coverage of 90·1\%. No
single predictor was strongly discriminatory in isolation (maximum
univariable AUC 0·65), indicating that performance derived from
multivariable integration rather than any dominant marker. Performance
was heterogeneous across subgroups: discrimination was consistently
lower in patients aged \(\ge\)65 years, and the modest negative predictive
value (62·0\%) at the operating threshold, alongside reduced sensitivity
in the older stratum, identified clear limits of a routine-data-only
approach in these groups.

\textbf{Interpretation} Using only routine pre-procedural data,
CardioPredict discriminated obstructive CAD with good calibration and
quantified case-level uncertainty, supporting its potential role as a
decision-support aid for pre-angiographic triage. Its novel contribution
resides not in a single biomarker but in three methodological features
combined in one pipeline: leakage-free repeated nested cross-validation,
clinical-grade calibration, and a conformal abstention mechanism that
explicitly flags statistically indeterminate patients rather than forcing
a binary decision. However, discrimination was weaker in patients aged
\(\ge\)65 years, the negative predictive value at the operating threshold was
modest (62\%), and a substantial minority (23.7\%) fell in the gray-zone;
because the model was developed in a single, angiography-referred cohort
with high disease prevalence, external and prospective validation, and
recalibration to local prevalence, are required before clinical
deployment.

\textbf{Funding} {[}Funding source / None{]}.

\subsection{Research in context}\label{research-in-context}

\paragraph{Evidence before this study}\label{evidence-before-this-study}

Pre-test probability tools for obstructive CAD---such as the updated
Diamond--Forrester and CAD Consortium models---rely chiefly on age, sex
and symptom typology and tend to over-estimate disease in contemporary
populations, particularly in lower-risk and female patients. Recent guidance
continues to endorse risk-factor-weighted likelihood models, and
calibrated risk-factor and calcium-score-weighted likelihood models have
recently been developed against CT-defined obstructive disease\textsuperscript{17,22}, and
machine-learning estimators have been shown to improve discrimination over
these classical scores in both invasive and CT-referred cohorts\textsuperscript{18,19,20,21}.
However, even the most recent machine-learning pre-test models report
modest discrimination when restricted to routine data alone (AUC
\(\approx\)0.77--0.80)\textsuperscript{18,19,20}, are frequently limited by
optimistic single-split validation, inadequate calibration assessment, and
undeclared label leakage from procedure-derived variables, and generally
fail to quantify case-level uncertainty or abstain from a decision for
genuinely ambiguous patients\textsuperscript{20,21}.

\paragraph{Added value of this study}\label{added-value-of-this-study}

We report an ensemble model estimated under repeated nested
cross-validation with explicit, audited removal of angiography- and
outcome-derived predictors, formal calibration assessment, and conformal
case-level uncertainty that explicitly abstains on ambiguous patients.
Performance reflects honest out-of-sample estimates rather than a single
optimistic split. In contrast to recent machine-learning pre-test models
that restrict discrimination to \(\approx\)0.77--0.80 with routine clinical
variables\textsuperscript{18,19,20}, CardioPredict achieved an AUC of 0.866 using
the same class of routinely available data, combined with good
calibration (E/O ratio 1.02) and an interpretable abstention mechanism
that reserves a statistically justified gray-zone for 23.7\% of patients
instead of forcing a binary decision. These design features are intended
directly to address the methodological limitations that have limited the
clinical translation of earlier models.

\paragraph{Implications of all the available
evidence}\label{implications-of-all-the-available-evidence}

A calibrated, uncertainty-aware estimate of obstructive-CAD probability
from routinely available data could help prioritise invasive angiography
and reduce avoidable procedures. Multi-centre external validation,
prospective evaluation, and prevalence recalibration are prerequisites
for clinical use.

\subsection{Introduction}\label{introduction}

Coronary artery disease (CAD) remains the leading cause of morbidity and mortality worldwide, accounting for over nine million deaths annually.\textsuperscript{1} Invasive coronary angiography (ICA) is the diagnostic reference standard for obstructive disease and frequently the gateway to revascularisation. However, ICA is resource-intensive, exposes patients to procedural and radiation risks, and remains associated with a high rate of normal or non-obstructive findings. Large-scale registry data indicate that up to 60\% of elective diagnostic catheterisations reveal no obstructive stenosis, representing a substantial diagnostic yield deficit and unnecessary risk exposure.\textsuperscript{2} Accurate estimation of the pre-test probability (PTP) of obstructive CAD from information available before catheterisation is therefore central to rational patient selection and to reducing low-value invasive testing.

For decades, clinical guidelines have relied on risk stratification scores to guide diagnostic decision-making. The landmark Diamond--Forrester model (1979) established the paradigm of using age, sex, and chest pain typicality to estimate the pre-test likelihood of CAD.\textsuperscript{3} This framework was subsequently updated and refined by the CAD Consortium models, which incorporated basic risk factors such as diabetes, hypertension, smoking, and dyslipidaemia.\textsuperscript{4} Despite their widespread adoption in guidelines, such as those from the European Society of Cardiology (ESC),\textsuperscript{5} these classical models are known to overestimate the probability of obstructive CAD in contemporary cohorts. Large-scale validation studies, including those based on the CONFIRM registry, have demonstrated that traditional scores overestimate disease prevalence by up to two-to-three-fold, particularly in women and lower-risk populations.\textsuperscript{6} This overestimation is primarily attributed to secular declines in cardiovascular risk factors and the widespread use of primary preventive therapies---such as statins and antihypertensives---over the last forty years, which have shifted the baseline pre-test likelihood compared to the 1970s and 1980s development cohorts.

To overcome the limitations of simple demographic-based scores, machine learning (ML) has emerged as a promising tool to integrate high-dimensional clinical inputs. Modern machine-learning architectures can capture complex, non-linear interactions among patient characteristics, blood biomarkers (including advanced lipid sub-fractions, glycaemic indicators, and systemic inflammatory markers), and detailed clinical histories. Studies leveraging large multicentre databases, such as the CONFIRM registry and the PROMISE trial, have demonstrated that ML algorithms---such as gradient-boosted decision trees and ensemble methods---significantly outperform traditional clinical scores in predicting both obstructive CAD and long-term major adverse cardiovascular events.\textsuperscript{7,8,9} This pattern has persisted in the most recent literature: contemporary machine-learning pre-test probability models, including the C-STRAT score in a large Chinese registry\textsuperscript{18} and multimodal deep-learning estimators combining clinical features with electrocardiographic information\textsuperscript{19,20}, consistently outperform classical symptom-based scores, although their discrimination with routine clinical data alone tends to plateau at an AUC of approximately 0.77--0.80\textsuperscript{18,20,21}. By incorporating laboratory measurements and biomarkers already routinely collected in secondary care, these models present a scalable solution for pre-angiographic triage without requiring additional imaging or invasive testing; nevertheless, the incremental value of pure multimodal signal (e.g.\ electrocardiographic or calcium-scoring data) over routine variables has been variable in external validation\textsuperscript{19,20,21}.

However, the clinical translation and utility of machine learning in this domain have been severely undermined by recurrent methodological limitations in the literature. First, many published models suffer from \emph{label leakage} or data contamination, incorporating predictors that are only determined during or after the angiographic procedure (such as procedure-related complications or post-test medication changes), which artificially inflates model performance. Second, models are frequently validated on a single optimistic data split, or using cross-validation schemes that fail to isolate preprocessing steps (such as imputation, outlier winsorisation, or feature engineering) within the training loop, leading to information leakage from the test set and overoptimistic performance. Third, the majority of research reports discrimination metrics (such as the Area Under the Receiver Operating Characteristic curve, AUC) while omitting formal assessments of probability calibration, leaving clinicians with risk scores whose absolute values cannot be trusted for individual patient triage. Finally, standard classifiers force a binary classification on all cases, disregarding the inherent clinical uncertainty of borderline patients and leading to potentially dangerous misclassifications.

We set out to develop and internally validate a clinical decision-support model---CardioPredict---that estimates the probability of obstructive CAD using only pre-procedural routine data, with four explicit design priorities that directly address these shortcomings. First, we establish an audited predictor taxonomy of 30 variables to prevent any label leakage. Second, we employ repeated nested cross-validation (two repeats of five-fold nested cross-validation) to report honest, leakage-free, out-of-sample performance estimates in accordance with the TRIPOD statement.\textsuperscript{10} Third, we implement age-stratified probability calibration to deliver reliable, clinical-grade probabilities. Fourth, we integrate split-conformal prediction to quantify case-level uncertainty, establishing a mathematically sound diagnostic ``gray-zone'' (abstention mechanism) that routes uncertain cases to non-invasive imaging while providing high-confidence decisive predictions for the remaining cohort.\textsuperscript{11} We benchmark this ensemble of tabular foundation models (TabPFN),\textsuperscript{12} gradient-boosted trees (LightGBM\textsuperscript{13} and CatBoost\textsuperscript{14}), and clinical logistic regression against standard models to validate its clinical utility. This combination of leakage-free estimation, clinical-grade calibration and a statistically principled abstention mechanism---applied jointly within a single routine-data pipeline---constitutes the study's principal innovation relative both to classical scores and to recent machine-learning pre-test models that still impose unconditional binary predictions and single-split validation\textsuperscript{18,20,21}.

\subsection{Methods}\label{methods}

\subsubsection{Study design and
participants}\label{study-design-and-participants}

We performed a single-centre retrospective analysis of patients referred
for clinically indicated coronary angiography. The source registry
contained 2346 entries. Patients without a recorded angiographic outcome
(1159 of 2346; 49·4\%) were ineligible and were excluded, yielding an
analytic sample of 1187 patients. Patients with and without a recorded
angiographic outcome differed substantially on fully-recorded baseline
variables (appendix Table A1): those who proceeded to angiography were
older and had markedly higher fasting glucose, triglycerides,
high-sensitivity C-reactive protein and blood pressure, and lower total
and LDL cholesterol --- a profile consistent with selective referral of
higher-risk patients. This verification/spectrum bias is a key
limitation (see Discussion) and means the cohort is not representative
of an unselected chest-pain population. No duplicate records were
detected using a seven-variable composite key (age, sex, fasting
glucose, total cholesterol, body-mass index, triglycerides,
haemoglobin).

\subsubsection{Outcome}\label{outcome}

The prediction target was angiographically defined obstructive coronary
artery disease (CAD), defined as ≥50\% luminal diameter stenosis in at
least one major epicardial coronary artery (or ≥50\% left main stenosis)
on invasive coronary angiography, recorded as a binary outcome. Controls
(no obstructive CAD) were patients referred for and undergoing
angiography in whom no epicardial stenosis reached this threshold ---
i.e. angiographically verified, not a separate unscreened "healthy"
sample --- so the outcome is defined identically in cases and controls.
The analytic sample comprised 781 patients with obstructive CAD (65·8\%)
and 406 without (34·2\%); the class-imbalance ratio was 1·92. This
prevalence reflects an angiography-referred population and is
substantially higher than in screening or general-outpatient settings.

\subsubsection{Predictors and prevention of label
leakage}\label{predictors-and-prevention-of-label-leakage}

From 223 raw fields we removed, by a prespecified taxonomy, identifier
and free-text columns (n=5), all angiography/vessel-level fields and
other outcome-derived variables (n=69), and zero-variance columns (n=8),
leaving 141 candidate variables that were used for the descriptive
cohort audit. From these we used a \emph{prespecified} set of \textbf{30
routine pre-procedural predictors} as model inputs (this fixed list, not
a data-driven selection, is given in full in the appendix). Three
variables with \textgreater50\% missingness were deliberately excluded
from the primary model --- platelet count (71\% missing;
non-discriminatory, group means 155·5 vs 154·6, p=0·86) and the
depression and anxiety scores (\textasciitilde52\% missing) --- to avoid
imputation-driven artefacts (a sensitivity analysis confirmed
performance was unchanged by their inclusion). The 30 predictors span
demographics (age, sex), anthropometry (body-mass index, waist
circumference, waist-to-hip ratio), blood pressure and hypertension
status, glycaemic and insulin-resistance indices (fasting glucose,
diabetes, triglyceride--glucose index), the lipid profile and derived
atherogenic ratios (total, LDL and HDL cholesterol, non-HDL,
total-cholesterol/HDL ratio, atherogenic index of plasma, and a
composite triglyceride·hs-CRP·glucose/HDL index), systemic inflammation
and haematology (high-sensitivity C-reactive protein, white-cell count,
haemoglobin), lifestyle (smoking history, physical-activity level),
medication (statin use), and personal/family history (prior myocardial
infarction, hypertension, and family history of diabetes, cardiovascular
disease and hypertension). Within each cross-validation training fold
these predictors were expanded with clinically motivated, deterministic
engineered features (haemodynamic composites, atherogenic and
insulin-resistance indices, and guideline-based thresholds), giving 76
modelled features (events-per-variable ≈ 10, with regularised base
learners); engineering used no outcome information.

\subsubsection{Missing data and outlier
handling}\label{missing-data-and-outlier-handling}

Structurally absent smoking and alcohol sub-fields were treated as
missing-not-at-random and set to zero (n=9 fields). Remaining missing
values were imputed with the median; to avoid leakage of distributional
information across folds, imputation statistics were estimated within
each cross-validation training partition and applied to the held-out
partition, and on the full cohort for the deployed model.
Physiologically implausible values were winsorised to fixed clinical
bounds (constants, independent of the data), and composite indices were
recomputed from the winsorised constituents. No outcome information
entered any preprocessing step.

\subsubsection{Model development}\label{model-development}

CardioPredict is a stacked ensemble of four complementary base learners:
a tabular foundation model (TabPFN),\textsuperscript{12} two gradient-boosted decision-tree
models (LightGBM\textsuperscript{13} and CatBoost\textsuperscript{14}), and a regularised clinical
logistic-regression model. Base-learner predictions were combined by an
age-conditional meta-learner (stacked generalisation), allowing the
optimal weighting of evidence to vary across age strata. Hyperparameters
for the gradient-boosted learners were tuned once on the full cohort by
age-stratified cross-validation using Bayesian (tree-structured Parzen)
optimisation (LightGBM, 150 trials; CatBoost, 75 trials) and held fixed
thereafter (Figure 1). We treat this single full-cohort tuning as a
fixed model-specification step that is held constant across the
performance-estimation folds (rather than re-searched within each inner
loop); because the search is target-aware, it can introduce mild
optimism, which we flag as a limitation. Given the modest class
imbalance (65·8\% events) we did not apply resampling or class
reweighting; the operating point is instead managed through calibration
and explicit threshold selection, preserving probability
interpretability.

\begin{figure}
\centering
\includegraphics[width=0.82\linewidth,height=\textheight,keepaspectratio,alt={Model pipeline schematic}]{outputs/figures/fig1_pipeline_schematic}
\caption{CardioPredict ensemble pipeline --- from routine data through
fold-internal preprocessing and the four base learners to the
age-conditional meta-learner, calibration, and conformal prediction.}
\end{figure}

\subsubsection{Calibration and conformal
uncertainty}\label{calibration-and-conformal-uncertainty}

Predicted probabilities were calibrated using an age-stratified
calibrator fitted on out-of-fold predictions. Case-level uncertainty was
quantified by split-conformal prediction\textsuperscript{11} at a 90\% target coverage: each
patient receives a prediction set, and cases whose calibrated
probability falls within the indeterminate "gray-zone" (where both
obstructive and non-obstructive labels remain statistically plausible)
are flagged for further evaluation rather than dichotomised.

\subsubsection{Decision thresholds}\label{decision-thresholds}

We derived several operating points from out-of-fold predictions to
support different clinical goals: a primary threshold maximising the
Matthews correlation coefficient (0·73, which also coincided with the
Youden-optimal point); an age-stratified threshold (0·50 for patients
aged \textless65 years; 0·73 for those aged ≥65 years); and a
high-sensitivity, cost-sensitive rule-out threshold (0·18) penalising
false negatives.

\subsubsection{Statistical analysis}\label{statistical-analysis}

Performance was estimated by 2 repeats of 5-fold nested
cross-validation. In every outer fold, imputation, scaling, feature
engineering, the four base learners, the meta-learner and the calibrator
were re-fitted on the training partition only and applied to the
untouched test partition; pooled out-of-fold predictions over all
patients yielded the reported metrics, summarised as the mean and
standard deviation across repeats. Discrimination was assessed by AUC;
calibration by the calibration slope, calibration-in-the-large and the
expected/observed (E/O) event ratio; and clinical utility by
sensitivity, specificity, predictive values, accuracy and the Matthews
correlation coefficient at the operating threshold. Univariable
discrimination of each predictor was computed to characterise individual
signal strength. The final deployed model was fitted on 100\% of the
cohort, with thresholds, calibrator and the conformal quantile derived
from leakage-free full-cohort out-of-fold predictions, and all fitted
base learners persisted to guarantee deterministic, reproducible
inference. Ninety-five-percent confidence intervals for the AUC and
operating-point metrics were obtained by bootstrap resampling (2000
replicates) of the pooled out-of-fold predictions. To quantify
incremental value we benchmarked the ensemble against a penalised
logistic-regression model fitted under the identical folds and
fold-internal preprocessing, comparing AUCs by paired bootstrap;
established symptom-based scores (Diamond--Forrester\textsuperscript{3}; CAD Consortium\textsuperscript{4})
were not computable because chest-pain typology was not recorded.
Conformal performance was summarised by the empirical coverage of the
prediction sets and the proportion of patients assigned to the
indeterminate (gray-zone) set. Analyses used a fixed random seed (42);
software versions are reported in the appendix. Reporting follows the
Transparent Reporting of a multivariable prediction model for Individual
Prognosis Or Diagnosis (TRIPOD) statement.\textsuperscript{10}

\subsubsection{Role of the funding
source}\label{role-of-the-funding-source}

{[}The funder had no role in study design, data collection, analysis,
interpretation, or writing of the report. / Not applicable.{]} The
corresponding author had full access to all data and final
responsibility for the decision to submit.

\subsection{Results}\label{results}

Of 2346 registry entries, 1187 patients (50·6\%) had an angiographic
outcome and constituted the analytic sample; 781 (65·8\%) had
obstructive CAD. Patients with obstructive CAD were older and had higher
fasting glucose, triglycerides, high-sensitivity C-reactive protein,
white-cell count and systolic and diastolic blood pressure, and lower
HDL cholesterol, than those without (Table 1). Total and LDL
cholesterol, body-mass index and platelet count did not differ
significantly between groups (Figure 2).

\begin{figure}
\centering
\includegraphics[width=0.72\linewidth,height=\textheight,keepaspectratio,alt={Study flow diagram}]{outputs/figures/fig2_study_flow}
\caption{Study profile. Derivation of the analytic cohort from the
angiography registry, with exclusions and class breakdown.}
\end{figure}

\begin{longtable}[]{@{}lllll@{}}
\caption{Table 1: Baseline characteristics of the analytic cohort, by
angiographic outcome (N=1187)}\tabularnewline
\toprule\noalign{}
Variable & Unit & No obstructive CAD (n=406) & Obstructive CAD (n=781) &
p value* \\
\midrule\noalign{}
\endfirsthead
\toprule\noalign{}
Variable & Unit & No obstructive CAD (n=406) & Obstructive CAD (n=781) &
p value* \\
\midrule\noalign{}
\endhead
\bottomrule\noalign{}
\endlastfoot
Age & years & 53·7 (11·4) & 58·9 (10·6) & \textless0·001 \\
Body-mass index & kg/m² & 26·9 (5·2) & 27·1 (5·2) & 0·47 \\
Waist circumference & cm & 93·7 (12·8) & 94·4 (12·5) & 0·35 \\
Systolic blood pressure & mm Hg & 129·8 (23·7) & 134·5 (25·0) & 0·004 \\
Diastolic blood pressure & mm Hg & 79·9 (12·0) & 81·9 (11·2) & 0·006 \\
Fasting blood glucose & mg/dL & 115·7 (46·8) & 133·1 (62·6) &
\textless0·001 \\
Total cholesterol & mg/dL & 168·2 (43·7) & 169·1 (42·9) & 0·71 \\
Triglycerides & mg/dL & 138·1 (71·3) & 150·0 (73·4) & \textless0·001 \\
HDL cholesterol & mg/dL & 42·5 (11·6) & 40·8 (12·1) & \textless0·001 \\
LDL cholesterol & mg/dL & 97·7 (35·2) & 98·2 (34·5) & 0·82 \\
High-sensitivity CRP & mg/L & 5·3 (6·6) & 6·9 (8·8) & \textless0·001 \\
White-cell count & ×10³/µL & 7·0 (2·3) & 7·7 (3·8) & \textless0·001 \\
Haemoglobin & g/dL & 13·0 (1·6) & 13·4 (1·7) & \textless0·001 \\
Platelet count & ×10³/µL & 155·5 (68·1) & 154·6 (59·1) & 0·86 \\
\end{longtable}

Data are mean (SD). *Mann--Whitney U test. CAD=coronary artery disease.
CRP=C-reactive protein.

\begin{figure}
\centering
\includegraphics[width=0.72\linewidth,height=\textheight,keepaspectratio,alt={Baseline standardised mean differences}]{outputs/figures/fig3_baseline_smd}
\caption{Baseline biomarker differences by angiographic outcome,
expressed as standardised mean differences (Cohen\textquotesingle s d;
positive = higher in CAD).}
\end{figure}

\begin{figure}
\centering
\includegraphics[width=0.64\linewidth,height=\textheight,keepaspectratio,alt={Missing-data profile}]{outputs/figures/fig4_missingness}
\caption{Proportion of missing values (before imputation) for each model
predictor; the dashed line marks 10\%. Imputation was fitted within each
cross-validation fold.}
\end{figure}

\subsubsection{Individual predictor
signal}\label{individual-predictor-signal}

No single predictor was strongly discriminatory: the highest univariable
AUCs were for diabetes (0·65), sex (0·65), a composite
triglyceride·hs-CRP·glucose/HDL metabolic-inflammatory index (0·64), and
age (0·63). The contrast between weak univariable signal and the
multivariable AUC of 0·87 indicates that discrimination arose from the
integration of many modestly informative variables rather than from any
dominant marker (Figure 5).

\begin{figure}
\centering
\includegraphics[width=0.72\linewidth,height=\textheight,keepaspectratio,alt={Univariate AUC ranking}]{outputs/figures/fig5_univariate_auc}
\caption{Top predictors ranked by univariable (single-variable)
discrimination; no single variable exceeds an AUC of
\textasciitilde0·65.}
\end{figure}

\subsubsection{Discrimination and
calibration}\label{discrimination-and-calibration}

Under repeated nested cross-validation with fold-internal imputation the
model achieved an AUC of 0·866 (SD 0·001 across repeats; per-repeat
values 0·864 and 0·867; bootstrap 95\% CI 0·850--0·891). The narrow
between-repeat SD reflects the large per-fold sample and the stability
of the calibrated ensemble; because only two repeats were run, the
bootstrap interval---rather than the across-repeat SD---should be read
as the primary measure of precision. Calibration was good, with a
calibration slope of 0·91, calibration-in-the-large of −0·05 and an E/O
ratio of 1·02 (Table 2; Figure 6, Figure 7).

\begin{longtable}[]{@{}lll@{}}
\caption{Table 2: Cross-validated performance of CardioPredict at the
primary operating threshold (0·73), with 95\% bootstrap confidence
intervals}\tabularnewline
\toprule\noalign{}
Metric & Value & 95\% CI \\
\midrule\noalign{}
\endfirsthead
\toprule\noalign{}
Metric & Value & 95\% CI \\
\midrule\noalign{}
\endhead
\bottomrule\noalign{}
\endlastfoot
Area under the ROC curve (AUC) & 0·866 & 0·85--0·89 \\
Sensitivity & 74·0\% & 70·8--77·0 \\
Specificity & 81·5\% & 77·6--85·2 \\
Positive predictive value & 88·5\% & 86·0--91·0 \\
Negative predictive value & 62·0\% & 58·0--66·0 \\
Overall accuracy & 76·6\% & 74·1--79·0 \\
Matthews correlation coefficient & 0·53 & --- \\
Calibration slope & 0·91 & --- \\
Calibration-in-the-large & −0·05 & --- \\
Expected/observed (E/O) ratio & 1·02 & --- \\
Conformal coverage (target / empirical) & 90\% / 90·1\% & --- \\
\end{longtable}

Predictive values are reported at the development-cohort prevalence
(65·8\%) and should be recalibrated to the local setting.

\begin{figure}
\centering
\includegraphics[width=0.58\linewidth,height=\textheight,keepaspectratio,alt={ROC curve}]{outputs/figures/fig6_roc}
\caption{Discrimination. Pooled nested cross-validation ROC curve (bold)
with the two per-repeat curves (faint); area under the curve 0·866 (SD
0·001).}
\end{figure}

\begin{figure}
\centering
\includegraphics[width=0.58\linewidth,height=\textheight,keepaspectratio,alt={Calibration plot}]{outputs/figures/fig7_calibration}
\caption{Calibration. Observed event frequency versus predicted
probability by decile (slope 0·91; E/O 1·02), with the
predicted-probability distribution inset.}
\end{figure}

\subsubsection{Operating points and clinical
utility}\label{operating-points-and-clinical-utility}

At the primary (MCC-optimised) threshold the model favoured specificity
(81·5\%) over sensitivity (74·0\%), with a high positive predictive
value (88·5\%) --- an operating point suited to confidently ruling
disease \emph{in} and prioritising patients for angiography. This
confirmatory strength, however, comes at the cost of a modest negative
predictive value (62·0\%) and of missing roughly a quarter of cases
with obstructive disease; the model is therefore not suited, at this
threshold, to ruling disease out, and a low-probability result should
not override strong clinical suspicion. Because the
optimal cut-off depends on the clinical goal, we provide alternative
operating points (Table 3): a high-sensitivity, cost-sensitive threshold
(0·18) for rule-out/screening contexts, and an age-stratified threshold
to account for age-dependent prevalence. The threshold should be matched
to the intended use rather than applied as a single fixed value.

\begin{longtable}[]{@{}lll@{}}
\caption{Table 3: Decision thresholds derived from out-of-fold
predictions}\tabularnewline
\toprule\noalign{}
Operating point & Probability cut-off & Intended use \\
\midrule\noalign{}
\endfirsthead
\toprule\noalign{}
Operating point & Probability cut-off & Intended use \\
\midrule\noalign{}
\endhead
\bottomrule\noalign{}
\endlastfoot
Primary (MCC-optimised; = Youden) & 0·73 & Balanced; rule-in /
prioritise angiography \\
Age \textless65 years (MCC) & 0·50 & Age-adjusted assessment \\
Age ≥65 years & 0·73 & Older patients \\
Cost-sensitive (high-sensitivity) & 0·18 & Rule-out / screening \\
\end{longtable}

\begin{figure}
\centering
\includegraphics[width=0.58\linewidth,height=\textheight,keepaspectratio,alt={Confusion matrix}]{outputs/figures/fig8_confusion}
\caption{Classification at the primary threshold (0·73), with
sensitivity, specificity, and predictive values.}
\end{figure}

\begin{figure}
\centering
\includegraphics[width=0.58\linewidth,height=\textheight,keepaspectratio,alt={Decision-curve analysis}]{outputs/figures/fig9_decision_curve}
\caption{Decision-curve analysis. Net benefit of CardioPredict versus
treat-all and treat-none strategies across threshold probabilities.}
\end{figure}

\subsubsection{Incremental value over a simple
model}\label{incremental-value-over-a-simple-model}

To quantify the added value of the ensemble we fitted a penalised
logistic-regression model under the identical 2×5 cross-validation folds
with fold-internal preprocessing. Given the \emph{same engineered
features} as the ensemble, the logistic model reached an AUC of 0·794
(95\% CI 0·77--0·82); the ensemble improved discrimination by +0·077
(95\% CI +0·059 to +0·098; paired bootstrap p\textless0·001), a
statistically significant gain. Critically, feature engineering alone
explained little of this difference --- a logistic model on the raw
predictors scored 0·789, so the engineered features added only ≈0·005,
and the +0·077 advantage is attributable to the non-linear ensemble
rather than to feature construction. Established symptom-based scores
established symptom-based scores (Diamond--Forrester\textsuperscript{3}, CAD Consortium\textsuperscript{4}) could not be reproduced because
chest-pain typology was not recorded; the logistic model therefore
serves as the routine-data reference. Among the individual base learners
(descriptive full-cohort out-of-fold AUC), TabPFN was strongest (0·874),
followed by LightGBM (0·867) and CatBoost (0·860), with the clinical
logistic learner lowest (0·778); the calibrated ensemble matched the
best single learner while adding calibration and uncertainty
quantification. Predictor-importance analysis (SHAP)\textsuperscript{16} of a representative
gradient-boosted base learner identified high-sensitivity C-reactive
protein, fasting glucose, age, sex (and sex-by-lipid interactions), and
triglycerides as the dominant contributors (Figure 10), consistent with
an inflammatory--metabolic risk signature.

\begin{figure}
\centering
\includegraphics[width=0.72\linewidth,height=\textheight,keepaspectratio,alt={Feature importance}]{outputs/figures/fig10_feature_importance}
\caption{Predictor importance (relative gain) for a representative
gradient-boosted base learner, illustrating which routine variables most
influence the ensemble\textquotesingle s predictions.}
\end{figure}

\subsubsection{Uncertainty
quantification}\label{uncertainty-quantification}

Applying the conformal procedure at a 90\% target coverage, 23·7\% of
patients fell within the indeterminate gray-zone (calibrated probability
between 0·32 and 0·68), where both obstructive and non-obstructive
labels remained statistically plausible; the remaining 76·3\% received a
decisive prediction. Empirical coverage of the prediction sets was
90·1\%, meeting the nominal target. In practice, a gray-zone result
identifies the patient in whom the routine-data model is least
informative and for whom further evaluation (e.g. non-invasive imaging)
is most likely to change management, whereas decisive predictions can
support a more confident triage decision (Figure 12).

\subsubsection{Subgroup performance}\label{subgroup-performance}

Discrimination tended to be higher in younger patients and lower in
older patients (≥65 years), in whom routine-data discrimination is
intrinsically harder; in this older stratum the model's discrimination
and sensitivity were visibly attenuated, indicating that a
routine-data-only signal offers the least leverage exactly where
pretest probability is highest. Disease prevalence differed markedly by
sex in this cohort, and subgroup specificity and sensitivity varied
accordingly; positive predictions in women and negative predictions in
men therefore warrant additional scrutiny (Figure 11). This asymmetric,
subgroup-dependent behaviour reinforces the case for careful threshold
selection and for caution when applying a single fixed operating point
across heterogeneous populations. The conformal framework assigned
statistically indeterminate cases to a probability
"gray-zone", separating them from decisive predictions (Figure 12).

\begin{figure}
\centering
\includegraphics[width=0.58\linewidth,height=\textheight,keepaspectratio,alt={Subgroup AUC plot}]{outputs/figures/fig11_subgroup_auc}
\caption{Discrimination by subgroup. Nested cross-validation AUC (95\%
bootstrap CI) overall and by sex and age band.}
\end{figure}

\begin{figure}
\centering
\includegraphics[width=0.72\linewidth,height=\textheight,keepaspectratio,alt={Conformal gray-zone}]{outputs/figures/fig12_conformal_grayzone}
\caption{Conformal uncertainty. Distribution of calibrated probabilities
by true class, with the shaded indeterminate gray-zone (90\% coverage)
in which both labels remain plausible.}
\end{figure}

\subsection{Discussion}\label{discussion}

In a cohort of 1187 patients referred for coronary angiography, a
calibrated machine-learning ensemble estimated the probability of
obstructive CAD from routine pre-procedural data with an AUC of 0·86 and
good calibration, and provided explicit case-level uncertainty. Because
the model uses only variables available before catheterisation, it is
positioned as a decision-support aid for pre-angiographic triage rather
than a replacement for invasive assessment.

Three features distinguish this work methodologically. First,
performance was estimated by repeated nested cross-validation with all
preprocessing and learners re-fitted within each training fold, avoiding
the optimism of single-split validation. Second, predictors derived from
the angiographic procedure or the outcome were removed by an audited
taxonomy, eliminating a common and under-recognised source of inflated
performance. Third, conformal prediction supplies an interpretable
abstention mechanism: rather than forcing every patient into a binary
label, statistically ambiguous cases are flagged for further evaluation,
which is clinically more faithful than a single dichotomous output.

The multivariable AUC of 0·87, compared with a maximum univariable AUC
of 0·65, indicates that clinical value derives from integrating multiple
modestly informative variables rather than any single biomarker. The
ensemble improved discrimination by a significant +0·077 over a
transparent logistic-regression model given the same engineered
features; because a logistic model on the raw variables performed almost
identically (0·789 vs 0·794), this gain reflects the non-linear ensemble
itself rather than feature construction. Good calibration (slope 0·91,
E/O 1·02) is important for a decision-support tool, since clinicians
must be able to trust the magnitude of the estimated probability, not
only its rank ordering.

The model showed positive net benefit over both treat-all and treat-none
strategies across a clinically relevant range of threshold probabilities
(Figure 9), supporting its use to guide referral.\textsuperscript{15} We set the primary
threshold to maximise the Matthews correlation coefficient because it
balances all four confusion-matrix quantities under class imbalance; on
this model it coincided with the Youden-optimal point (0·73) and
favoured specificity (81·5\%) and positive predictive value (88·5\%)
over sensitivity (74·0\%) --- an operating point suited to confirmatory
triage and prioritising angiography. For rule-out or screening contexts
a lower, high-sensitivity threshold (0·18) is provided, and we report
several operating points (Table 3) so the cut-off can be matched to the
local goal rather than applied as a single fixed value. The conformal
gray-zone (23·7\% of patients, empirical coverage 90·1\%)
operationalises diagnostic uncertainty: these are precisely the patients
in whom routine data are insufficient and further testing is most likely
to change management, whereas the 76·3\% with decisive predictions can
be triaged with greater confidence. Discrimination was lower in patients
aged ≥65 years, so the model should be applied with extra caution in
older patients, in whom non-invasive imaging may remain preferable.

Interpreted alongside the most recent machine-learning pre-test work,
these results also define clear boundaries of the proposed approach.
While our ensemble improved on the ≈0.77--0.80 ceiling reported for
routine-data-only estimators in several recent cohorts\textsuperscript{18,20,21},
the absolute margin was modest relative to multimodal models that add
electrocardiographic or calcium-scoring information\textsuperscript{19,20}, and all
subsequent comparisons rely on different, non-overlapping populations
and endpoints. In particular, models incorporating coronary-artery
calcium scores consistently report substantially higher discrimination
(AUC ≈0.87--0.90)\textsuperscript{20}, underscoring that routine laboratory data
alone---however carefully combined---cannot match the information
contained in direct anatomical imaging. Our design therefore does not
position CardioPredict as a substitute for such imaging but as a
zero-cost triage step that reserves it for the patients most likely to
benefit.

\subsubsection{Limitations}\label{limitations}

This study has important limitations. It is a single-centre,
retrospective analysis of an angiography-referred population with high
disease prevalence (65·8\%); the resulting spectrum and verification
bias mean that the predictive values reported here will not transfer
directly to lower-prevalence screening or general-outpatient settings
without recalibration. The model has not yet undergone external or
geographic validation, nor prospective evaluation, both of which are
essential before clinical use. Performance varied by sex and age
subgroup, reflecting differing prevalence and limited subgroup sample
sizes, and was intrinsically weaker in older patients (≥65 years) and
in whom routine-data signal is least informative. The sample size
(1187), while adequate for development, constrains precision for rare
subgroups. At the primary threshold the modest negative predictive value
(62·0\%) and sensitivity (74·0\%) mean the model is best reserved for
confirmatory triage rather than rule-out, and the sizable gray-zone
(23·7\%) indicates that routine data were insufficient for a decisive
prediction in a substantial minority of patients---an honest ceiling of
the data modality rather than a model shortcoming. Finally, the
predictors are limited to routine data; incorporating information with
established incremental value for obstructive CAD---such as
coronary-artery-calcium scoring, cardiac troponin, electrocardiographic
or stress-test findings---would be expected to raise discrimination
beyond the present ceiling\textsuperscript{19,20}.

\subsubsection{Positioning within the clinical
pathway}\label{positioning-within-the-clinical-pathway}

In current practice, a patient with suspected CAD undergoes clinical
assessment and non-invasive testing (resting or stress
electrocardiography, stress imaging, or coronary CT angiography); those
with positive or inconclusive results are referred for invasive coronary
angiography, and only a subset are found to have obstructive disease
warranting percutaneous coronary intervention (PCI). Because referral
rests on tests of limited specificity, a substantial proportion of
diagnostic angiograms reveal no obstructive stenosis --- an avoidable
exposure to procedural risk and cost (Figure 13, current route).

CardioPredict inserts a single, calibrated, uncertainty-aware triage
step immediately after routine clinical and laboratory data are obtained
--- before any invasive procedure --- using information already
collected at no additional cost (Figure 13, proposed route). A decisive
high-probability output (high specificity and positive predictive value
at the operating threshold) identifies patients who can be prioritised
for angiography and likely PCI; a decisive low-probability output
supports deferral to medical management, reducing unnecessary
catheterisation; and the conformal gray-zone explicitly routes genuinely
uncertain patients to targeted non-invasive testing rather than directly
to the catheter laboratory. Importantly, the tool is designed to augment
--- not replace --- existing tests and clinical judgement, and (pending
external validation and recalibration to local prevalence) a
low-probability result should not override strong clinical suspicion.

\begin{figure}
\centering
\includegraphics[width=0.72\linewidth,height=\textheight,keepaspectratio,alt={Clinical-pathway diagram}]{outputs/figures/fig13_clinical_pathway}
\caption{CardioPredict in the diagnostic pathway. From a shared clinical
assessment, current practice uses non-invasive testing to refer positive
or inconclusive patients to invasive coronary angiography. The proposed
triage instead converts routine pre-procedural data into a calibrated
probability with a conformal gray-zone, sending decisive
high-probability patients to angiography, deferring decisive
low-probability patients to medical management, and routing gray-zone
patients to the same non-invasive testing. Both routes converge on the
same invasive endpoint, so the tool changes \emph{which} patients reach
the catheter laboratory, not the procedures that follow. Navy =
CardioPredict model; red = invasive angiography; green = medical
management; dashed box = proposed approach.}
\end{figure}

\subsubsection{Implications and future
work}\label{implications-and-future-work}

Subject to external and prospective validation, a calibrated,
uncertainty-aware estimate of obstructive-CAD probability could help
cardiologists prioritise invasive angiography and reduce avoidable
diagnostic procedures, with potential benefits for patient safety and
resource use. Priorities for further research are multi-centre external
validation with prevalence recalibration, prospective impact evaluation
on referral decisions, and assessment of incremental value when
established imaging and biomarker variables are added.

\subsection{Declarations}\label{declarations}

\subsubsection{Contributors}\label{contributors}

{[}Initials{]} conceived the study and obtained funding. {[}Initials{]}
curated the data and implemented the analysis pipeline. {[}Initials{]}
and {[}Initials{]} verified the underlying data and analyses.
{[}Initials{]} drafted the manuscript. All authors contributed to
interpretation, critically revised the manuscript, approved the final
version, and had final responsibility for the decision to submit.

\subsubsection{Declaration of interests}\label{declaration-of-interests}

{[}We declare no competing interests. / Specify.{]}

\subsubsection{Data sharing}\label{data-sharing}

De-identified participant data, the analysis code, and the preprocessing
pipeline are available from the corresponding author on reasonable
request, subject to institutional and ethics-committee approval. The
fixed analysis seed and recorded software versions enable deterministic
reproduction.

\subsubsection{Ethics}\label{ethics}

{[}The study was approved by the {[}name{]} institutional review board
(reference {[}number{]}); the requirement for individual consent was
waived for this retrospective analysis of de-identified data. /
Specify.{]}

\subsubsection{Acknowledgements}\label{acknowledgements}

{[}Acknowledgements.{]}

\subsection{Appendix: reproducibility}\label{appendix-reproducibility}

Random seed 42 · development cohort N=1187 · data hash (SHA-256, 16 hex)
fb3af1d8645920f3 · Python 3.12 · NumPy 1.26.4 · scikit-learn 1.8.0 ·
LightGBM 4.6.0 · CatBoost 1.2.10 · TabPFN 7.1.1. Evaluation: 2×5
repeated nested cross-validation; deployed model fitted on the full
cohort with persisted base learners for deterministic inference.

\begin{longtable}[]{@{}llll@{}}
\caption{Appendix Table A1: Baseline characteristics of included
(outcome recorded) versus excluded (outcome missing) registry
patients}\tabularnewline
\toprule\noalign{}
Variable & Included (n=1187) & Excluded (n=1159) & p value* \\
\midrule\noalign{}
\endfirsthead
\toprule\noalign{}
Variable & Included (n=1187) & Excluded (n=1159) & p value* \\
\midrule\noalign{}
\endhead
\bottomrule\noalign{}
\endlastfoot
Age, years & 57·1 (11·2) & 53·0 (9·3) & \textless0·001 \\
Male, \% & 48·6 & 49·6 & 0·63 \\
Fasting glucose, mg/dL & 127·2 (58·3) & 84·0 (17·8) & \textless0·001 \\
Total cholesterol, mg/dL & 168·7 (43·3) & 182·8 (34·3) &
\textless0·001 \\
Triglycerides, mg/dL & 145·9 (72·9) & 114·7 (59·5) & \textless0·001 \\
HDL cholesterol, mg/dL & 41·6 (14·2) & 44·5 (9·5) & \textless0·001 \\
LDL cholesterol, mg/dL & 98·0 (34·8) & 113·3 (30·8) & \textless0·001 \\
High-sensitivity CRP, mg/L & 6·4 (8·2) & 1·4 (0·7) & \textless0·001 \\
Body-mass index, kg/m² & 27·0 (5·2) & 26·8 (4·2) & 0·49 \\
Systolic BP, mmHg & 132·9 (24·7) & 120·6 (15·5) & \textless0·001 \\
Diastolic BP, mmHg & 81·3 (11·5) & 74·9 (9·9) & \textless0·001 \\
\end{longtable}

Data are mean (SD) unless stated. *Mann--Whitney U test (sex: χ²).
Included patients (who underwent angiography) carried a substantially
higher cardiometabolic and inflammatory burden than excluded patients,
confirming selective referral; predictive values derived in this cohort
therefore require recalibration before use in lower-risk populations.

\subsection{References}\label{references}

\begin{enumerate}
\tightlist
\item
  Roth GA, Mensah GA, Johnson CO, et al. Global burden of cardiovascular
  diseases and risk factors, 1990--2019. \emph{J Am Coll Cardiol}.
  2020;76(25):2982--3021.
\item
  Patel MR, Peterson ED, Dai D, et al. Low diagnostic yield of elective coronary angiography. \emph{N Engl J Med}. 2010;362(10):886--95.
\item
  Diamond GA, Forrester JS. Analysis of probability as an aid in the
  clinical diagnosis of coronary-artery disease. \emph{N Engl J Med}.
  1979;300(24):1350--8.
\item
  Genders TSS, Steyerberg EW, Hunink MGM, et al. Prediction model to
  estimate presence of coronary artery disease: retrospective pooled
  analysis of existing cohorts. \emph{BMJ}. 2012;344:e3485.
\item
  Knuuti J, Wijns W, Saraste A, et al. 2019 ESC Guidelines for the diagnosis and management of chronic coronary syndromes. \emph{Eur Heart J}. 2020;41(3):407--77.
\item
  Cheng VY, Carter AM, et al. Performance of the traditional age, sex, and angina typicality-based approach for estimating pretest probability of angiographically significant coronary artery disease in patients undergoing coronary computed tomographic angiography. \emph{Circulation}. 2011;124(22):2423--32.
\item
  Motwani M, Dey D, et al. Machine learning for prediction of all-cause mortality in patients with suspected coronary artery disease: a 5-year multicentre prospective registry analysis. \emph{Eur Heart J}. 2017;38(7):500--7.
\item
  Al'Aref SJ, Maliakal G, et al. Machine learning of clinical variables and coronary artery calcium scoring for the prediction of obstructive coronary artery disease on coronary computed tomography angiography: analysis from the CONFIRM registry. \emph{Eur Heart J}. 2020;41(3):359--67.
\item
  Douglas PS, Hoffmann U, Patel MR, et al. Outcomes of anatomical versus functional testing for coronary artery disease. \emph{N Engl J Med}. 2015;372(14):1291--300.
\item
  Collins GS, Reitsma JB, Altman DG, Moons KGM. Transparent Reporting of
  a multivariable prediction model for Individual Prognosis Or Diagnosis
  (TRIPOD): the TRIPOD statement. \emph{BMJ}. 2015;350:g7594.
\item
  Vovk V, Gammerman A, Shafer G. \emph{Algorithmic Learning in a Random
  World}. New York: Springer; 2005.
\item
  Hollmann N, Müller S, Purucker L, et al. Accurate predictions on small
  data with a tabular foundation model. \emph{Nature}.
  2025;637(8045):319--26.
\item
  Ke G, Meng Q, Finley T, et al. LightGBM: a highly efficient gradient
  boosting decision tree. \emph{Adv Neural Inf Process Syst (NeurIPS)}.
  2017;30:3146--54.
\item
  Prokhorenkova L, Gusev G, Vorobev A, Dorogush AV, Gulin A. CatBoost:
  unbiased boosting with categorical features. \emph{Adv Neural Inf
  Process Syst (NeurIPS)}. 2018;31:6638--48.
\item
  Vickers AJ, Elkin EB. Decision curve analysis: a novel method for
  evaluating prediction models. \emph{Med Decis Making}.
  2006;26(6):565--74.
\item
  Lundberg SM, Lee S-I. A unified approach to interpreting model
  predictions. \emph{Adv Neural Inf Process Syst (NeurIPS)}.
  2017;30:4765--74.
\item
  Vrints C, Andreotti F, Koskinas KC, et al. 2024 ESC Guidelines for the
  management of chronic coronary syndromes. \emph{Eur Heart J}.
  2024;45(36):3415--3537.
\item
  Dou G, Zhou J, Guo Z, et al. Developing a machine-learning-based
  pre-test probability model of obstructive coronary artery disease in
  patients with stable chest pain (C-STRAT). \emph{JACC Asia}.
  2025; doi:10.1016/j.jacasi.2025.03.015.
\item
  Trivedi RK, Chiu I-M, Hughes JW, Rogers AJ, Ouyang D. Deep learning on
  electrocardiogram waveforms to stratify risk of obstructive stable
  coronary artery disease. \emph{Eur Heart J Digit Health}. 2025;
  doi:10.1093/ehjdh/ztaf020.
\item
  Chen T, Shao D, Zhao J, et al. Comparison of the RF-CL and CACS-CL
  models to estimate the pretest probability of obstructive coronary
  artery disease and predict prognosis in patients with stable chest
  pain and diabetes mellitus. \emph{Front Cardiovasc Med}. 2024;11:1368743.
\item
  Hennecken J, Verpalen VA, van de Vijver WR, et al. Optimization of
  pre-test probability models for obstructive coronary artery disease
  using explainable artificial intelligence. \emph{Eur Heart J}.
  2025; doi:10.1093/eurheartj/ehaf784.186.
\item
  Rasmussen LD, Schmidt SE, Knuuti J, et al. Clinical likelihood models
  calibrated against observed obstructive coronary artery disease on
  computed tomography angiography. \emph{Eur Heart J Cardiovasc Imaging}.
  2025; doi:10.1093/ehjci/jeaf043.
\end{enumerate}

\end{document}
